\paragrafos{Ιδιότητες της ισότητας και η έννοια της εξίσωσης}{Α.6.1, «Ιδιότητες της ισότητας (ισότητα και πράξεις) και η έννοια της εξίσωσης»}

\begin{anagke}
Στο εργαστήριο του σχολείου ο Ηλίας βρίσκει μια παλιά ζυγαριά με δύο δίσκους. Στον αριστερό δίσκο βάζει ένα κλειστό σακουλάκι με φουντούκια, που δεν ξέρει πόσο ζυγίζει, και ένα βαρίδι των $100$ g. Στον δεξί δίσκο βάζει βαρίδια που ζυγίζουν συνολικά $350$ g. Η ζυγαριά ισορροπεί. «Πόσο ζυγίζουν τα φουντούκια;» ρωτάει η Ρήνη.

Ο Στέλιος αρπάζει το βαρίδι των $100$ g από τον αριστερό δίσκο. Η ζυγαριά γέρνει δεξιά. «Δεν πειράζει», λέει η Φανή, «αν αφαιρέσουμε $100$ g και από τον άλλο δίσκο, θα ξαναισορροπήσει». Έτσι κάνουν. Στον αριστερό μένει το σακουλάκι και στον δεξιό $250$ g. Η Πολυξένη το γράφει στο τετράδιό της και το κάνει ισότητα: «Πρόσθεσα ή αφαίρεσα το ίδιο και από τις δύο πλευρές, και η ισότητα έμεινε αληθινή». Ο Βασίλης αμφιβάλλει: «Και αν ήταν πολλαπλασιασμός; Ή διαίρεση;».

Στην ισότητα αυτή υπάρχει ένας αριθμός που δεν ξέρουμε, το βάρος του σακουλιού. Σε αυτή την παράγραφο θα μάθουμε ποιες πράξεις μπορούμε να κάνουμε χωρίς να χαλάσουμε μια ισότητα, και θα γνωρίσουμε την εξίσωση, δηλαδή την ισότητα που κρύβει έναν άγνωστο αριθμό.
\end{anagke}

\begin{lexi}{ισότητα, εξίσωση, άγνωστος, ρίζα, επαλήθευση}{ισότητα, εξίσωση, ρίζα}
Η λέξη \lexiword{εξίσωση} βγαίνει από το «ίσος»: εξισώνω σημαίνει «κάνω ίσα» τα δύο μέλη, όπως οι δύο δίσκοι μιας ζυγαριάς. \lexiword{Άγνωστος} είναι ο αριθμός που δεν ξέρουμε και ψάχνουμε. Η \lexiword{ρίζα} της εξίσωσης είναι η λύση της, και λέγεται έτσι γιατί «βγαίνει από το βάθος», όπως η ρίζα ενός φυτού. Η \lexiword{επαλήθευση} είναι ο έλεγχος ότι κάτι είναι αληθινό («επί» + «αληθεύω»).
\end{lexi}

\begin{metafrasi}
Η ζυγαριά ισορροπεί. & $\alpha=\beta$\\
Βάζω το ίδιο βάρος και στους δύο δίσκους. & $\alpha+\gamma=\beta+\gamma$\\
Βγάζω το ίδιο βάρος και από τους δύο δίσκους. & $\alpha-\gamma=\beta-\gamma$\\
Διπλασιάζω ή τριπλασιάζω ό,τι υπάρχει και στους δύο δίσκους. & $\alpha\cdot\gamma=\beta\cdot\gamma$\\
Χωρίζω και τους δύο δίσκους σε ίσα μέρη. & $\alpha:\gamma=\beta:\gamma$, με $\gamma\neq0$\\
Σακουλάκι και $100$ g ισορροπούν με $350$ g. & $x+100=350$\\
\end{metafrasi}

\tmima{Η θεωρία}

\tmimasmall{Ισότητα και ισορροπία}

Μια ισότητα $\alpha=\beta$ λέει ότι οι δύο πλευρές της, που λέγονται \textbf{μέλη}, έχουν την ίδια τιμή. Η ζυγαριά που ισορροπεί είναι η καλύτερη εικόνα για αυτό.

\begin{center}
\scalebox{0.78}{\binzyg{\binsx\ \binsw{100 g}}{\binsw{350 g}}}
\quad\raisebox{6mm}{$\Rightarrow$}\quad
\scalebox{0.78}{\binzyg{\binsx}{\binsw{250 g}}}
\end{center}

Αν βγάλουμε $100$ g από τον έναν δίσκο και $100$ g από τον άλλον, η ζυγαριά ισορροπεί ακόμη και φαίνεται ότι $x=250$ g.

\begin{idiotita}{ισότητα και πράξεις}
Για οποιουσδήποτε αριθμούς $\alpha,\beta,\gamma$ ισχύει:
\begin{itemize}[leftmargin=1.2em,itemsep=0.3ex,topsep=2pt]
\item Αν $\alpha=\beta$, τότε $\alpha+\gamma=\beta+\gamma$ (προσθέτουμε τον ίδιο αριθμό και στα δύο μέλη).
\item Αν $\alpha=\beta$, τότε $\alpha-\gamma=\beta-\gamma$ (αφαιρούμε τον ίδιο αριθμό και από τα δύο μέλη).
\item Αν $\alpha=\beta$, τότε $\alpha\cdot\gamma=\beta\cdot\gamma$ (πολλαπλασιάζουμε και τα δύο μέλη με τον ίδιο αριθμό).
\item Αν $\alpha=\beta$, τότε $\alpha:\gamma=\beta:\gamma$ (διαιρούμε και τα δύο μέλη με τον ίδιο αριθμό \textbf{που δεν είναι μηδέν}).
\end{itemize}
\end{idiotita}

Για παράδειγμα, από την $8=8$ προκύπτουν οι $8+5=8+5$, $8\cdot3=8\cdot3$ και $8:2=8:2$, δηλαδή $13=13$, $24=24$ και $4=4$. Οι ιδιότητες αυτές είναι το εργαλείο μας για όλη την Άλγεβρα.

\begin{prosoxi}
Δεν διαιρούμε με το $0$. Από την $0\cdot3=0\cdot5$ (που ισχύει, γιατί $0=0$) δεν μπορούμε να «διαιρέσουμε με το $0$» και να καταλήξουμε στο $3=5$. Για αυτό στην τελευταία ιδιότητα απαιτείται $\gamma\neq0$.
\end{prosoxi}

\tmimasmall{Η εξίσωση}

\begin{orismos}{εξίσωση με έναν άγνωστο}
\textbf{Εξίσωση} με έναν άγνωστο λέγεται μια ισότητα που περιέχει αριθμούς και ένα γράμμα (τον άγνωστο), π.χ. $x+5=15$ ή $2\cdot y=16$.\\
\textbf{Λύση ή ρίζα} της εξίσωσης είναι ο αριθμός που, όταν αντικαταστήσει τον άγνωστο, επαληθεύει την ισότητα. Η διαδικασία εύρεσης της λύσης λέγεται \textbf{επίλυση} της εξίσωσης.
\end{orismos}

Ο αριθμός $10$ είναι λύση της $x+5=15$, γιατί όταν αντικαταστήσουμε το $x$ με $10$ προκύπτει $10+5=15$, που ισχύει. Ο αριθμός $7$ δεν είναι λύση, γιατί $7+5=12\neq15$. Μετά από κάθε επίλυση κάνουμε \textbf{επαλήθευση}: βάζουμε τη λύση στην \emph{αρχική} εξίσωση και ελέγχουμε ότι βγαίνει σωστή ισότητα.

\tmimasmall{Ειδικές περιπτώσεις}

\begin{itemize}[leftmargin=1.2em,itemsep=0.3ex,topsep=2pt]
\item Η εξίσωση $0\cdot x=2$ δεν επαληθεύεται από κανέναν αριθμό, γιατί κάθε αριθμός επί $0$ κάνει $0$. Μια εξίσωση χωρίς καμία λύση λέγεται \textbf{αδύνατη}.
\item Η εξίσωση $0\cdot x=0$ επαληθεύεται από όλους τους αριθμούς. Μια εξίσωση με λύσεις όλους τους αριθμούς λέγεται \textbf{ταυτότητα} ή \textbf{αόριστη}.
\end{itemize}

\tmimasmall{Ισοδύναμες εξισώσεις και οι έξι βασικές μορφές}

Δύο εξισώσεις που έχουν τις ίδιες λύσεις λέγονται \textbf{ισοδύναμες}. Οι ιδιότητες της ισότητας μας επιτρέπουν να αντικαθιστούμε μια εξίσωση με μια ισοδύναμη, πιο απλή. Από τον ορισμό των πράξεων (η αφαίρεση είναι η αντίστροφη της πρόσθεσης και η διαίρεση του πολλαπλασιασμού) βρίσκουμε τη λύση κάθε μίας από τις παρακάτω εξισώσεις.

\begin{center}
\renewcommand{\arraystretch}{1.5}
\begin{tabular}{@{}lll@{}}
\toprule
\textbf{Εξίσωση} & \textbf{Λύση} & \textbf{Παράδειγμα}\\
\midrule
$x+\alpha=\beta$ & $x=\beta-\alpha$ & $x+2=5\ \Rightarrow\ x=5-2=3$\\
$x-\alpha=\beta$ & $x=\beta+\alpha$ & $x-6=2\ \Rightarrow\ x=2+6=8$\\
$\alpha-x=\beta$ & $x=\alpha-\beta$ & $10-x=3\ \Rightarrow\ x=10-3=7$\\
$\alpha\cdot x=\beta$ & $x=\beta:\alpha$ & $2x=40\ \Rightarrow\ x=40:2=20$\\
$x:\alpha=\beta$ & $x=\beta\cdot\alpha$ & $x:5=6\ \Rightarrow\ x=6\cdot5=30$\\
$\alpha:x=\beta$ & $x=\alpha:\beta$ & $20:x=5\ \Rightarrow\ x=20:5=4$\\
\bottomrule
\end{tabular}
\end{center}

(Στις εξισώσεις με διαίρεση υποθέτουμε ότι δεν διαιρούμε με το $0$.) Οι αριθμοί των εξισώσεων μπορεί να είναι και ρητοί, με πρόσημο, δεκαδικοί ή κλάσματα.

\begin{prosoxi}
\begin{itemize}[leftmargin=1.2em,itemsep=0.3ex,topsep=2pt]
\item Στην $\alpha-x=\beta$ η λύση είναι $\alpha-\beta$, όχι $\beta-\alpha$ ή $\alpha+\beta$. Η $12-x=5$ δίνει $x=7$, όχι $17$.
\item Στις $\alpha:x=\beta$ και $x:\alpha=\beta$ οι λύσεις διαφέρουν: $24:x=6$ δίνει $x=24:6=4$, ενώ $x:6=24$ δίνει $x=24\cdot6=144$.
\item Στην επαλήθευση βάζουμε τη λύση στην \emph{αρχική} εξίσωση, όχι σε κάποια που προέκυψε στη μέση της λύσης μας.
\end{itemize}
\end{prosoxi}

\tmima{Λυμένα παραδείγματα}

\paradeigma{Η ζυγαριά του Ηλία}
\tixero Σακουλάκι με $x$ g φουντούκια και βαρίδι $100$ g ισορροπούν με $350$ g.\par
\tithelo Το βάρος των φουντουκιών.\par
\possindeo Η ισότητα είναι $x+100=350$. Αφαιρώ $100$ και από τα δύο μέλη (ιδιότητα της αφαίρεσης): $x+100-100=350-100$, δηλαδή $x=250$.\par
\apantisi $250$ g.\par
\elegxos $250+100=350$. Σωστό.

\paradeigma{Είναι ο αριθμός λύση;}
\tixero Η εξίσωση $2x=16$ και οι αριθμοί $6$ και $8$. Η εξίσωση $x-4=1$ και οι αριθμοί $3$ και $5$.\par
\tithelo Ποιοι από αυτούς είναι ρίζες.\par
\possindeo Αντικαθιστώ. Για $x=6$: $2\cdot6=12\neq16$. Για $x=8$: $2\cdot8=16$ σωστό. Για $x=3$: $3-4=-1\neq1$. Για $x=5$: $5-4=1$ σωστό.\par
\apantisi Ρίζες είναι ο $8$ και ο $5$.\par
\elegxos Από τους κανόνες: $x=16:2=8$ και $x=1+4=5$.

\paradeigma{Λύνω τις έξι βασικές εξισώσεις}
\tixero Οι εξισώσεις $x+2{,}5=7$, $x-\frac{3}{4}=\frac{1}{4}$, $15-x=6{,}4$, $4x=18$, $x:6=2{,}5$, $36:x=9$.\par
\tithelo Να τις λύσω.\par
\possindeo $x=7-2{,}5=4{,}5$. $x=\frac{1}{4}+\frac{3}{4}=1$. $x=15-6{,}4=8{,}6$. $x=18:4=4{,}5$. $x=2{,}5\cdot6=15$. $x=36:9=4$.\par
\apantisi $4{,}5$, $1$, $8{,}6$, $4{,}5$, $15$, $4$.\par
\elegxos $4{,}5+2{,}5=7$, $1-\frac{3}{4}=\frac{1}{4}$, $15-8{,}6=6{,}4$, $4\cdot4{,}5=18$, $15:6=2{,}5$, $36:4=9$. Όλα σωστά.

\paradeigma{Ρητοί αριθμοί με πρόσημο}
\tixero Οι εξισώσεις $-5x=25$ και $x+(-8)=-12$ και $-8-x=11$.\par
\tithelo Να τις λύσω και να επαληθεύσω.\par
\possindeo Στην πρώτη διαιρώ και τα δύο μέλη με $-5$: $x=25:(-5)=-5$. Στη δεύτερη προσθέτω $8$: $x=-12+8=-4$. Στην τρίτη, αν προσθέσω $8$ και στα δύο μέλη, παίρνω $-x=19$, άρα $x=-19$.\par
\apantisi $-5$, $-4$, $-19$.\par
\elegxos $-5\cdot(-5)=25$. $-4+(-8)=-12$. $-8-(-19)=-8+19=11$. Σωστά.

\paradeigma{Ταυτότητα ή αδύνατη;}
\tixero Οι εξισώσεις $x+3=3+x$, $x+2=x+6$ και $5x-5x=7$.\par
\tithelo Τι είδους είναι.\par
\possindeo Η πρώτη ισχύει για κάθε αριθμό (αντιμεταθετική ιδιότητα): ταυτότητα. Στη δεύτερη το αριστερό μέλος είναι πάντα κατά $4$ μικρότερο από το δεξί: αδύνατη. Στην τρίτη το $5x-5x=0$ και δεν ισούται με $7$: αδύνατη.\par
\apantisi Ταυτότητα, αδύνατη, αδύνατη.\par
\elegxos Με $x=10$: $13=13$, αλλά $12\neq16$ και $0\neq7$.

\paradeigma{Η δεξαμενή του χωριού}
\tixero Μια δεξαμενή έχει όγκο $7{,}2$ m$^3$, μήκος $2$ μ. και πλάτος $1{,}2$ μ.\par
\tithelo Το ύψος της.\par
\possindeo Αν $x$ είναι το ύψος σε μέτρα, ο όγκος είναι $2\cdot1{,}2\cdot x=2{,}4x$. Άρα $2{,}4x=7{,}2$ και $x=7{,}2:2{,}4=3$.\par
\apantisi $3$ μ.\par
\elegxos $2\cdot1{,}2\cdot3=7{,}2$. Σωστό.

\tmima{Ασκήσεις}

\zoni{Το κατάλαβα}{ιδιότητες, ρίζες, βασικές εξισώσεις}

\ex Σωστό ή λάθος; α) Αν $\alpha=\beta$, τότε $\alpha+5=\beta+5$. \quad β) Αν $x-3=7$, τότε $x=4$. \quad γ) Η εξίσωση $0\cdot x=3$ είναι αδύνατη. \quad δ) Αν $3\alpha=0$, τότε $\alpha=0$. \quad ε) Αν $\dfrac{\alpha}{3}=\dfrac{\beta}{3}$, τότε $\alpha=\beta$. \quad στ) Η εξίσωση $2x=6$ έχει λύση το $6$.

\ex \textbf{Η ζυγαριά.} Στην εικόνα η ζυγαριά ισορροπεί. Γράψε την εξίσωση και βρες πόσο ζυγίζει το σακουλάκι.
\begin{center}
\binzyg{\binsx\ \binsw{200 g}}{\binsw{500 g}}
\end{center}

\ex Αξιοποίησε τις ιδιότητες της ισότητας και συμπλήρωσε τα κενά. α) Αν $\alpha+2=7$, τότε $\alpha=\dots$ \quad β) Αν $\kappa-3=10$, τότε $\kappa=\dots$ \quad γ) Αν $3=\lambda$, τότε $\dots=3$ \quad δ) Αν $9\rho=45$, τότε $\rho=\dots$

\ex Εξέτασε αν οι αριθμοί είναι λύσεις των εξισώσεων: α) $x+3=9$, $x=6$ \quad β) $x-2=7$, $x=9$ \quad γ) $2x=10$, $x=20$ \quad δ) $12-x=10$, $x=22$ \quad ε) $x:4=3$, $x=12$.

\ex Λύσε τις εξισώσεις: α) $x+12=45$ \quad β) $9+y=17$ \quad γ) $x-48=5$ \quad δ) $65-x=17$ \quad ε) $15y=45$ \quad στ) $z:5=40$ \quad ζ) $26:x=13$.

\ex Λύσε τις εξισώσεις και επαλήθευσε: α) $-5x=25$ \quad β) $x+(-8)=-12$ \quad γ) $-6x=-18$ \quad δ) $-8-x=11$.

\ex Λύσε: α) $x+0{,}75=3$ \quad β) $\dfrac{1}{2}x=\dfrac{9}{4}$ \quad γ) $x-\dfrac{1}{3}=\dfrac{5}{6}$ \quad δ) $x:\dfrac{3}{5}=10$.

\zoni{Το χρησιμοποιώ}{ζυγαριές, μυστικοί αριθμοί, σχήματα}

\ex \textbf{Ο μυστικός αριθμός.} Η Πολυξένη λέει: «Σκέφτηκα έναν αριθμό». Γράψε εξίσωση και βρες τον, όταν: α) πρόσθεσε $17$ και βρήκε $52$ \quad β) τον πολλαπλασίασε επί $3$ και βρήκε $84$ \quad γ) τον αφαίρεσε από το $50$ και βρήκε $18$ \quad δ) τον διαίρεσε με $2$ και βρήκε $7$.

\ex \textbf{Τα βαρίδια.} Στη ζυγαριά της θείας Μαρίας δύο ίδια μπουκάλια ζυγίζουν όσο ένα βαρίδι $1\,200$ g. Γράψε την εξίσωση και βρες το βάρος κάθε μπουκαλιού.

\ex \textbf{Τα εισιτήρια και ο κουμπαράς.} α) Τέσσερα ίδια εισιτήρια πλοίου κοστίζουν $54$ €. Πόσο κάνει το καθένα; β) Ο Ηλίας είχε στον κουμπαρά $36$ €, ξόδεψε κάποιο ποσό και του έμειναν $11{,}5$ €. Πόσα ξόδεψε;

\ex \textbf{Οι περίμετροι.} α) Η Αγγελική φτιάχνει μια τετράγωνη κορνίζα με περίμετρο $98$ cm. Πόση είναι η πλευρά της; β) Ένα ισόπλευρο τρίγωνο έχει περίμετρο $4{,}5$ m. Πόση είναι η πλευρά του;

\ex Βρες το είδος κάθε εξίσωσης (ένας αριθμός λύση, αδύνατη ή ταυτότητα): α) $x+3=3+x$ \quad β) $0\cdot x=7$ \quad γ) $x+2=x+6$ \quad δ) $5x-5x=0$.

\zoni{Το σκέφτομαι}{γιατί δουλεύει, πού κάνουμε λάθος}

\ex Γιατί η $0\cdot x=0$ έχει άπειρες λύσεις και η $0\cdot x=5$ καμία; Εξήγησε με το «ποιος αριθμός επί $0$ δίνει...».

\ex Ο Στέλιος λύνει την $12-x=5$ και γράφει $x=12+5=17$. Λύνει και την $24:x=6$ και γράφει $x=24\cdot6$. Βρες τα λάθη του, γράψε τις σωστές λύσεις και επαλήθευσε.

\ex Γράψε μια εξίσωση με λύση το $6$, μια αδύνατη και μια ταυτότητα.

\ex Ποια ζεύγη εξισώσεων είναι ισοδύναμα; α) $x+3=7$ και $2x=8$ \quad β) $x-1=2$ και $x+1=2$ \quad γ) $x+1=x+2$ και $0\cdot x=5$. Εξήγησε.

\zoni{Για το σπίτι}{φτιάξε το δικό σου}

\ex Ο παππούς Γιάννης λέει: «Σκέφτηκα έναν αριθμό και το τριπλάσιό του είναι $2{,}4$». Ποιος είναι; Γράψε την εξίσωση.

\ex Φτιάξε μια καθημερινή ιστορία που να οδηγεί στην εξίσωση $x+12=30$ και ζήτησε από ένα μέλος της οικογένειάς σου να τη λύσει.
